\documentclass{beamer}
\usetheme{Madrid}
\usecolortheme{crane}
\usepackage{listings}
\usepackage{xcolor}

\definecolor{charcoal}{RGB}{44,44,44}
\definecolor{tuered}{RGB}{180,30,45}

\setbeamercolor{palette primary}{bg=charcoal,fg=white}
\setbeamercolor{palette secondary}{bg=tuered,fg=white}
\setbeamercolor{palette tertiary}{bg=charcoal,fg=white}
\setbeamercolor{palette quaternary}{bg=tuered,fg=white}
\setbeamercolor{structure}{fg=tuered}
\setbeamercolor{section in toc}{fg=charcoal}

\lstset{
  language=Java,
  basicstyle=\ttfamily\small,
  keywordstyle=\color{tuered}\bfseries,
  commentstyle=\color{gray},
  stringstyle=\color{teal},
  backgroundcolor=\color{black!5},
  frame=single,
  breaklines=true
}

\title{OO Thinking}
\subtitle{Week 12 — Tuesday | Chapter 10}
\author{Professor Pyatt}
\institute{COP 2250 — Java Programming | Pyatt Labs @ SCF}
\date{Spring 2026}

\begin{document}

\begin{frame}
  \titlepage
\end{frame}

% -------------------------------------------------------
\begin{frame}{Where We Are}
  \begin{block}{So far}
    \begin{itemize}
      \item Week 10: Designed a \texttt{Rectangle} class — fields, constructors, methods
      \item You understand \textbf{encapsulation}: private fields, public accessors
    \end{itemize}
  \end{block}
  \begin{block}{This week}
    \begin{itemize}
      \item Chapter 10: \textbf{OO Thinking}
      \item Immutable objects, instance vs.\ static, method overloading
      \item You will build a wrapper class: \texttt{MyString} (lab) and \texttt{MyInteger} (assignment)
    \end{itemize}
  \end{block}
\end{frame}

% -------------------------------------------------------
\begin{frame}[fragile]{Immutable Objects}
  \begin{lstlisting}
String s = "hello";
s.toUpperCase();
System.out.println(s);  // still "hello"
  \end{lstlisting}
  \begin{block}{Key Point}
    \texttt{toUpperCase()} returns a \textbf{new} String.\\
    It does \textbf{not} modify the original.
  \end{block}
  \begin{itemize}
    \item Immutable = safe to share
    \item The stored value never changes after construction
    \item Java's \texttt{String} class is immutable by design
  \end{itemize}
\end{frame}

% -------------------------------------------------------
\begin{frame}[fragile]{String vs StringBuilder}
  \begin{lstlisting}
// String: creates a new object every concatenation
String s = "hello";
s = s + " world";

// StringBuilder: modifies in place
StringBuilder sb = new StringBuilder("hello");
sb.append(" world");
System.out.println(sb.toString());
  \end{lstlisting}
  \begin{alertblock}{Performance}
    In a loop concatenating 1000 strings, \texttt{StringBuilder}\\
    is dramatically faster. \texttt{String} creates 1000 objects.
  \end{alertblock}
\end{frame}

% -------------------------------------------------------
\begin{frame}[fragile]{Instance vs Static Methods}
  \begin{lstlisting}
MyInteger m = new MyInteger(7);

// Instance method — needs an object
m.isEven();

// Static method — belongs to the class
MyInteger.isEven(4);
  \end{lstlisting}
  \begin{block}{Rule of Thumb}
    \begin{itemize}
      \item Instance: method needs the \textbf{stored value}
      \item Static: works on \textbf{any input}, no object required
    \end{itemize}
  \end{block}
\end{frame}

% -------------------------------------------------------
\begin{frame}[fragile]{Method Overloading}
  \begin{lstlisting}
// Three versions of isEven — same name, different params
public boolean isEven()              // instance
public static boolean isEven(int n)  // static, int
public static boolean isEven(MyInteger m) // static, MyInteger
  \end{lstlisting}
  \begin{block}{Java picks based on the argument type}
    Don't duplicate logic. Have the \texttt{MyInteger} version\\
    call the \texttt{int} version:
  \end{block}
  \begin{lstlisting}
public static boolean isEven(MyInteger m) {
    return isEven(m.getValue());  // reuse!
}
  \end{lstlisting}
\end{frame}

% -------------------------------------------------------
\begin{frame}[fragile]{Wrapper Class Design}
  \begin{block}{What is a wrapper class?}
    A class that stores a single primitive value and\\
    exposes methods to inspect or transform it.
  \end{block}
  \begin{lstlisting}
public class MyString {
    private String value;

    public MyString(String value) {
        this.value = value;
    }

    public String toUpperCase() {
        return value.toUpperCase(); // returns new — immutable
    }
}
  \end{lstlisting}
  \begin{itemize}
    \item \texttt{value} is private
    \item Methods return new values — they do not reassign \texttt{value}
  \end{itemize}
\end{frame}

% -------------------------------------------------------
\begin{frame}[fragile]{isPrime — The Logic}
  \begin{lstlisting}
public static boolean isPrime(int n) {
    if (n <= 1) return false;
    for (int i = 2; i <= Math.sqrt(n); i++) {
        if (n % i == 0) return false;
    }
    return true;
}
  \end{lstlisting}
  \begin{itemize}
    \item 1 is not prime
    \item Only check up to $\sqrt{n}$ — any larger factor has a smaller partner
    \item If no divisor found: prime
  \end{itemize}
\end{frame}

% -------------------------------------------------------
\begin{frame}[fragile]{parseInt Without Integer.parseInt()}
  \begin{lstlisting}
// "1234" → 1234
int result = 0;
for (int i = 0; i < s.length(); i++) {
    result = result * 10
           + Character.getNumericValue(s.charAt(i));
}
  \end{lstlisting}
  \begin{block}{The Algorithm}
    Shift left one decimal place each iteration.\\
    Multiply by 10, add the digit.
  \end{block}
  \begin{alertblock}{Rule}
    You cannot use \texttt{Integer.parseInt()} inside\\
    your own \texttt{parseInt()} — that is circular.
  \end{alertblock}
\end{frame}

% -------------------------------------------------------
\begin{frame}{Key Terms}
  \begin{description}
    \item[Wrapper class] Stores a single primitive; exposes methods
    \item[Immutable] Object whose state cannot change after construction
    \item[Instance method] Operates on \texttt{this} object's stored value
    \item[Static method] Belongs to the class; no object needed
    \item[Method overloading] Same name, different parameter types
    \item[\texttt{this}] Reference to the current object's field
  \end{description}
\end{frame}

% -------------------------------------------------------
\begin{frame}{Today's Lab — MyString}
  \begin{block}{Build a wrapper class for \texttt{String}}
    Implement in order:
    \begin{enumerate}
      \item \texttt{getValue()} — return stored value
      \item \texttt{length()} — character count
      \item \texttt{toUpperCase()} — return new uppercase string
      \item \texttt{reverse()} — loop, no \texttt{StringBuilder.reverse()}
      \item \texttt{isPalindrome()} — use your \texttt{reverse()}
      \item \texttt{contains(String sub)} — return true/false
    \end{enumerate}
  \end{block}
  \begin{alertblock}{Demo Requirement}
    Screen-share, compile live, explain one method in your own words.
  \end{alertblock}
\end{frame}

% -------------------------------------------------------
\begin{frame}{What's Next}
  \begin{block}{Assignment 10 — MyInteger (Liang 10.3)}
    Same pattern. Wrapper for \texttt{int}.\\
    Adds: overloaded static methods, \texttt{equals()}, \texttt{parseInt()}.
  \end{block}
  \begin{block}{Thursday}
    Game project kickoff.\\
    Stage 1: Game Design Document — no code, just concept.\\
    Your team HedgeDoc is already set up.
  \end{block}
  \vspace{0.5em}
  \centering\textit{Start simple. Get it running. Then make it better.}
\end{frame}

\end{document}
